% ---------------------------------------------------------------------------
% Planned additions. They are placed last so that they can be promoted into the main evaluation once the results exist.
\section{External evaluation (planned)}
\label{sec:external}

\todo{planned; not final}
The evaluation in \cref{sec:experiments} uses DENTEX only. This section describes additions that are planned
but not yet final: the official DENTEX test set, benchmarks on which competing dental systems report results, the
harness as a benchmark in its own right, a smaller sibling of the backbone trained with data parallelism, and a Cloud TPU
path. None of them changes the frozen protocol of \cref{sec:eval} or any number reported above.

\subsection{The official DENTEX test set}
\label{sec:official-test}

DENTEX also publishes an official test set of 250 images. We plan to report the final conditions of \cref{tab:main} on it,
once, in addition to the 50-image set used so far. That set is DENTEX's official \emph{validation} split, not its test
split, so choices made on it, such as selecting the grounding detector's weights, do not leak into the final test
numbers. The 250 test images are not used for training, development, or selection.

\subsection{Benchmarks used by competing systems}
\label{sec:external-benchmarks}

The closest dental systems report results on benchmarks that we did not build (\cref{sec:related}), so the most direct
comparison is to run the same benchmark. MMOral-OPG-Bench \citep{hao2025mmoral} is public. It contains 100 panoramic
radiographs with 491 closed-ended and 578 open-ended questions in categories such as teeth, pathology, and jaw, and it is
used by OralGPT-Plus, OralAgent, and OPGAgent. We plan to score three kinds of system with one LLM judge, using the
benchmark authors' prompts and scoring rules unchanged:
\begin{itemize}[leftmargin=*,itemsep=1pt,topsep=2pt]
  \item the raw Qwen3.5-9B backbone;
  \item the agent-assisted system: the trained agent first reads each radiograph and produces its findings (it never sees
  the question), and the findings are added as context for the same base model, which then answers the benchmark's
  official prompt. The paired difference to the raw backbone isolates what the agent adds, with model, prompt, and judge
  held fixed;
  \item the public OralGPT-Omni-7B and OralGPT-OPG-7B models, which were trained on benchmark-style data, whereas
  the agent-assisted system was not.
\end{itemize}
Scores that other papers report for the same model depend on the judge: OralGPT-Omni is reported at 45.31 in its own
paper \citep{hao2025oralgptomni} and at 59.06 in the table of OPGAgent \citep{yu2026opgagent}. We therefore re-score every
system ourselves and show published numbers only as separate rows with their source and judge. Every mean carries a
95\% interval from resampling radiographs, because the questions of one image are correlated, and a leakage screen checks
the benchmark images against the DENTEX training and validation images. Because the agent reports only the four
DENTEX diagnosis classes, we expect any gain to concentrate in the teeth and pathology categories rather than in
questions about past treatment, the jaw, or recommendations; this is a hypothesis, not a result. The other benchmarks are
out of reach for now: MMOral-Uni is gated and contains no panoramic images, and MMOral-X and OPG-Bench had no public
release when we checked, so we would ask the authors for access.

\subsection{The harness as a benchmark on Kaggle}
\label{sec:kaggle}

Tools, rules, and skills form a self-contained, versioned interface (\cref{sec:harness-serving}), so the harness can serve
as a benchmark itself: a set of difficult dental OPG tasks together with the tool registry, the rules, and the skills.
We plan to score our agent on these tasks and to compare it with the strongest frontier models available through the
Kaggle API, each run with and without the harness. Two quantities matter: where our agent stands among the frontier
models, and the uplift the harness gives each model in carrying out the tasks (\cref{tab:kaggle}).

\begin{table}[htbp]
  \centering
  \small
  \caption{Planned Kaggle comparison; every cell is filled from the benchmark run.}
  \label{tab:kaggle}
  \begin{tabular}{lccc}
    \toprule
    Model & Without harness & With harness & Uplift \\
    \midrule
    Our agent (Qwen3.5-9B, trained) & \todo{} & \todo{} & \todo{} \\
    Frontier model 1 \todo{name} & \todo{} & \todo{} & \todo{} \\
    Frontier model 2 \todo{name} & \todo{} & \todo{} & \todo{} \\
    Frontier model 3 \todo{name} & \todo{} & \todo{} & \todo{} \\
    \bottomrule
  \end{tabular}
\end{table}

\subsection{A smaller sibling with data parallelism}
\label{sec:sibling}

SFT and GRPO are the slow stages of the pipeline. The 4B variant of Qwen3.5 is small enough to be replicated on every
device, so training can use plain data parallelism instead of sharding one large model, which speeds up both stages. We
may move the main experiments to this size. If we do, the 9B and 4B models become an older, bigger sibling and a younger,
smaller one: the same recipe, data, and evaluation at two sizes, so that their paired comparison shows how much the
agentic recipe depends on model size.

\subsection{A Cloud TPU path}
\label{sec:tpu}

The same pipeline can also run on Cloud TPU. There the model sees the aspect-preserving views letterboxed onto fixed
canvases with a static number of vision slots, which avoids recompilation, whereas on GPU the views add no padding.
The hardware profile is stored with every result, and results from different profiles are never mixed. This path is not
yet part of the reported experiments.
